\documentclass[11pt,letterpaper]{article}
\usepackage[T1]{fontenc}
\usepackage{lmodern}
\usepackage[margin=0.85in,headheight=15pt]{geometry}
\usepackage{amsmath,amssymb,booktabs,tabularx,array,enumitem,xcolor,fancyhdr,microtype}
\usepackage[authoryear,round]{natbib}
\usepackage[colorlinks=true,linkcolor=teal,citecolor=teal,urlcolor=teal,breaklinks=true]{hyperref}
\definecolor{ink}{HTML}{173B3F}
\definecolor{teal}{HTML}{176C70}
\setlength{\parindent}{0pt}
\setlength{\parskip}{6pt}
\setlength{\emergencystretch}{2em}
\renewcommand{\arraystretch}{1.24}
\newcolumntype{L}[1]{>{\raggedright\arraybackslash}p{#1}}
\newcolumntype{Y}{>{\raggedright\arraybackslash}X}
\setlist{nosep,leftmargin=1.5em}
\pagestyle{fancy}\fancyhf{}
\fancyhead[L]{\small\textcolor{ink}{Complexification / companion reading}}
\fancyhead[R]{\small 18 November 2026}
\fancyfoot[L]{\small Conceptual research draft}\fancyfoot[R]{\thepage}
\newcommand{\topic}[2]{\clearpage\section{#1}\textit{#2}\par}
\newcommand{\status}[1]{\textcolor{teal}{\textbf{Status: #1.}}}
\newcommand{\takeaway}[1]{\par\smallskip\textcolor{ink}{\textbf{Takeaway.}} #1\par}
\newcommand{\Prb}{\mathbb{P}}
\newcommand{\E}{\mathbb{E}}
\newcommand{\indep}{\mathrel{\perp\!\!\!\perp}}
\newcommand{\TV}{d_{\mathrm{TV}}}
\newcommand{\references}{\clearpage\small\raggedright\bibliographystyle{plainnat}\bibliography{references}}

\title{What must people retain control over?\\\large Mathematical companion: reports, objectives, and decision rights}
\author{Sandy Tanwisuth}
\date{Research draft 3 / 18 November 2026}
\begin{document}\maketitle
\textbf{Core question.} Which information and decision rights let people correct AI-assisted collective decisions when circumstances or authorized objectives change?

This companion starts from strategic equivalence in \citet{lauffer2023knowledge}. It explains a simple refinement across decision contexts and an exactly checked example. This refinement is an elementary consequence of the definitions, not a claimed new theorem. The research proposal concerns learning and testing such interfaces under realistic limits on attention and coordination.

The AI-islands proposal associated with Alec's presentation is a motivation \citep{xor2026balance}. Human review of inter-assistant exchanges is a design to compare. The cited literature does not establish that universal human mediation is superior.

\begin{tabularx}{\linewidth}{L{0.25\linewidth}Y}\toprule
Status & What is available \\\midrule
Source foundation & Strategic equivalence; causal games; predictive feedback factorization; philosophical conditions for human control. \\
Checked here & An exhaustive three-co-policy illustration, including all five deterministic report partitions. \\
Proposed & Computational comparisons and a bounded organizational pilot. Neither has been run. \\
Unsettled & Which human reasons and objectives should govern; acceptable burden; numerical success thresholds and sample planning. \\\bottomrule
\end{tabularx}

\textbf{Reading route.} Section 1 fixes the notation. Sections 2--4 derive the reporting condition and work the example. Section 5 adds adaptation by other agents. Sections 6--7 connect this to an evaluable organizational design. Section 8 gives the source coverage; full citations follow.

\topic{Notation and four separate design choices}{A small information requirement can reflect a small action set.}
\begin{tabularx}{\linewidth}{L{0.25\linewidth}Y}\toprule
Symbol & Meaning \\\midrule
$i,\pi_i,\pi_{-i}$ & Focal decision maker, their policy, other participants' co-policy. \\
$\theta$ & Index for a possible co-policy; in the example $\theta\in\{C,D,E\}$. \\
$a\in\mathcal A_d$ & Available action under rights arrangement $d$. \\
$u_\lambda$ & Explicitly chosen payoff/objective indexed by $\lambda$. It is not inferred moral legitimacy. \\
$m=r(\theta)$ & Report of co-policy information in the toy model. A real report would be inferred from evidence, not oracle knowledge. \\
$\mathrm{BR}_{\lambda,d}(\theta)$ & Entire set of optimal feasible actions for that context. \\
$\mathcal Q$ & Declared family of contexts $q=(\lambda,d)$. \\
$V_q^{\rm full},V_q^r$ & Expected optimal value with full co-policy information, or with the report. \\
$K,O,r,d,\lambda$ & Generator; observation channel; report map; rights; objective. \\\bottomrule
\end{tabularx}

\textbf{Observation and reporting.} Masking a transcript changes $O$ or $r$ while a recorded generator stays fixed. It does not simulate removal of a communication channel.

\textbf{Generator.} Removing a channel or inserting human mediation changes $K$, including delays and possibly the policies of participants.

\textbf{Decision rights.} Changing $d$ changes available actions and who can trigger them. A permission on paper is not evidence that an intervention takes effect before the relevant deadline.

\textbf{Normative objective.} Changing $\lambda$ changes what counts as a better outcome. A more powerful collective need not serve each member, and agreement on a formal payoff does not establish legitimate authorization. Stakeholder conflict must be recorded rather than hidden in one scalar score.

\topic{Strategic equivalence, step by step}{Start with an existing result about coordination.}
For a finite decision problem, define
\begin{equation}
 \mathrm{BR}_i(\pi_{-i})=
 \arg\max_{\pi_i\in\Pi_i}R_i(\pi_i,\pi_{-i}).
\end{equation}
The symbol $\arg\max$ returns policies attaining the highest expected payoff, not the payoff itself. It can return more than one policy.

\textbf{Step 1: group co-policies that induce the same set.}
\begin{equation}
 \pi_{-i}\sim_i\pi'_{-i}\quad\Longleftrightarrow\quad
 \mathrm{BR}_i(\pi_{-i})=\mathrm{BR}_i(\pi'_{-i}).
\end{equation}
This is the strategic equivalence relation (SER) of \citet[Definition 4.3, PDF p. 4]{lauffer2023knowledge}. Its equivalence classes are the preimages of the best-response map. They are the largest cells that preserve the whole set unambiguously.

\textbf{Step 2: interpret ``minimal knowledge'' precisely.} If a report identifies the class, the player can recover their best-response set without distinguishing co-policies within it. This is a partition-based statement about information in a specified game. It is not a guarantee that a human can infer, read, trust, or act on the report.

\textbf{Step 3: distinguish a weaker requirement.} A cell $C$ can admit one action that is optimal throughout whenever
\begin{equation}
 \bigcap_{\theta\in C}\mathrm{BR}_i(\theta)\neq\varnothing.
\end{equation}
This does not require all the sets to be equal. For example, $\{A,B\}$ and $\{B,C\}$ share $B$ but differ. Such overlap need not define an equivalence relation, so it does not generally yield the same canonical partition as SER.

\textbf{Step 4: distinguish actions from values.} SER can merge co-policies with identical best responses but different achieved payoffs. The paper's valued SER additionally preserves the value of optimal responses. Its dynamic algorithm exploits recursive structure. Our one-step example below does not implement that algorithm.

The paper's computational treatment focuses on common-payoff games; the authors state that the underlying theory extends to general-sum games. We use a focal person's specified utility and do not assume that it exhausts the interests of all affected people.

\topic{Reports across a declared family of contexts}{The mathematical extension here is elementary; its usefulness is an open empirical question.}
For $q=(\lambda,d)$, define
\begin{equation}
 B_q(\theta)=\arg\max_{a\in\mathcal A_d}
 \E[u_\lambda\mid a,\theta].
\end{equation}
The expectation comes from the stipulated decision model. In a causal setting it must represent the consequences of choosing $a$, not merely conditioning on a correlated observed action.

\textbf{Step 1: decide what the report is meant to preserve.} If a fixed report must recover the best-response set for each context in $\mathcal Q$, it should satisfy
\begin{equation}\label{eq:report}
 r(\theta)=r(\theta')\ \Longrightarrow\
 B_q(\theta)=B_q(\theta')\quad\forall q\in\mathcal Q.
\end{equation}
The person is assumed to know which context $q$ applies. The report itself does not carry a new value judgment.

\textbf{Step 2: construct a sufficient report.} For a finite context family, let
\begin{equation}
 r_{\mathcal Q}(\theta)=(B_q(\theta))_{q\in\mathcal Q}.
\end{equation}
Its cells are intersections of the individual SER cells: the common refinement. If a report merges two types whose vectors differ, it fails equation~\eqref{eq:report} for at least one context. Thus these cells are the coarsest partition preserving all the listed sets. This short argument follows directly from the definition; no general novelty claim is made.

\textbf{Step 3: recognize the cost.} Adding contexts can only refine the partition or leave it unchanged. It cannot merge cells. More possible future corrections may require more distinctions now. But a system can instead issue a new report when the context changes, provided the underlying evidence remains accessible and the person has enough time. Fixed reports and adaptive reporting must be compared fairly.

\textbf{Research hypothesis.} Reports designed for a declared correction family may reduce regret after changed objectives relative to reports optimized only for routine decisions. Whether this works with uncertain evidence, learned co-policy models, strategic concealment, and human attention limits remains unknown.

\topic{Worked example: objective changes and action rights}{Every number below is an invented, unitless payoff.}
\status{checked illustration} Let $P(\theta=C)=P(\theta=D)=P(\theta=E)=1/3$ and $\mathcal A=\{A,B\}$. The three labels denote fixed collaborator co-policies. Only the focal objective changes in the first comparison.
\begin{center}\begin{tabular}{llrrr}\toprule
Objective & Action & C & D & E\\\midrule
Routine $u_0$ & A & 4&0&4\\
& B &2&3&2\\\midrule
Revised $u_1$ & A &4&0&0\\
& B &2&3&2\\\bottomrule\end{tabular}\end{center}
\textbf{Step 1: compute optimal actions column by column.} Under $u_0$, choose $A$ in C/E and $B$ in D. The SER is $\{\{C,E\},\{D\}\}$. Under $u_1$, choose $A$ in C and $B$ in D/E: $\{\{C\},\{D,E\}\}$. Preserving both requires three singleton cells.

\textbf{Step 2: evaluate the old report after the objective changes.} Full information gives
\begin{equation}
 V_1^{\rm full}=\tfrac13(4+3+2)=3.
\end{equation}
Inside the old C/E cell, each action has conditional expected value $2$. In D, $B$ gives $3$. Therefore
\begin{equation}
 V_1^{r_0}=\tfrac23(2)+\tfrac13(3)=\tfrac73,
 \qquad V_1^{\rm full}-V_1^{r_0}=\tfrac23.
\end{equation}
The loss is informational, with rights and co-policy prior fixed. If a new report can reveal C versus E in time, this loss can be avoided.

\textbf{Step 3: restrict rights separately.} Remove $B$, leaving $\mathcal A_d=\{A\}$. Every type now has the identical feasible best-response set $\{A\}$ under both objectives. A constant report preserves it. Yet
\begin{equation}
 V_{1,d}^{\rm full}=V_{1,d}^{\rm constant}=\tfrac13(4+0+0)=\tfrac43.
\end{equation}
The information regret under those restricted rights is zero, but attainable value has fallen. Compare rights arrangements separately; do not count the smaller information requirement as improved human control.

The script \citep{tanwisuth2026strategic} exhausts all five deterministic report partitions on three types and checks these values with rational arithmetic. This is not an empirical effect estimate or a general impossibility theorem.

\topic{Why adaptation by others changes the test}{A revised objective need not leave collaborators' policies fixed.}
The preceding illustration holds the co-policy prior and payoff table fixed except for the declared objective change. In a live organization, announcements can change what assistants and people choose to do. \citet[PDF pp. 15--17]{hammond2023games} distinguish:
\begin{itemize}
\item Post-policy intervention: alter an event while keeping the policy-selection outcome fixed.
\item Pre-policy intervention: make a change available when policies are selected, allowing strategic adaptation to the changed game.
\end{itemize}
This distinction concerns the information available to policy selection, not merely clock time. An announced permission or incentive change can alter other participants' responses. Those responses belong in a mechanized model, not in an unchanged empirical conditional distribution.

\citet{jorgensen2026collective} asks when a collective can be represented as a goal-directed macro agent under aligned interventions. Even if such a surrogate predicts the collective accurately, it does not establish that an affected person can alter its objective or prevent an action. Agency of the collective and agency of its members can diverge.

\textbf{Connection to world modeling.} \citet{rosas2025vat} provides controlled interfaces over possible actions; predicting one routine trajectory is weaker. The neural feedback paper \citep[PDF pp. 4--9, 24--25]{ehrlich2026feedback} explains how hiding interactions can disrupt a predictor's factorization. That result is about observation and representation. Removing the interaction is a generator change, and adding a human veto is a rights change.

\textbf{Why dimension is not the objective.} A report can discard predictable distinctions while retaining all actions needed for a task. Another report can preserve detailed predictions while concealing a reason to revise the task itself. The relevant measure therefore depends on what people are authorized to correct and when the evidence and actions are available.

\takeaway{Use replay to isolate information loss; use live interaction to study adaptation and mediation. Do not treat the former as evidence for the latter.}

\topic{From the equations to an AI-islands comparison}{The reporting map, communication graph, rights, and objective need separate interventions.}
The proposed task is fictional resource allocation. A handoff contains a recommendation, evidence references, uncertainty, assumptions, and dissent. The person reviewing it can query source evidence and revise or refuse the handoff. Final execution rights and an appeal route are specified independently.

\begin{tabularx}{\linewidth}{L{0.26\linewidth}Y}\toprule
Design & Main manipulation \\\midrule
Logged collaboration & Direct assistant exchanges with provenance and reviewable reports. \\
Human-mediated islands & Human approval and forwarding of cross-team exchanges. \\
Delegation with escalation & Routine exchanges continue; specified uncertainty or conflict triggers review. \\\bottomrule
\end{tabularx}

\textbf{Stage A: information comparison.} Keep the generator and rights fixed. Compare reports in routine cases and declared changed contexts. In a finite oracle task,
\begin{align}
 V_q^{\rm full}&=\sum_\theta p(\theta)\max_{a\in\mathcal A_d}u_q(a,\theta),\\
 V_q^r&=\sum_m p(m)\max_{a\in\mathcal A_d}
                  \sum_{\theta:r(\theta)=m}p(\theta\mid m)u_q(a,\theta),\\
 \mathrm{Regret}_q(r)&=V_q^{\rm full}-V_q^r\geq0.
\end{align}
The nonnegativity follows because a full-information decision maker can imitate the report-based policy. This benchmark assumes known payoffs, prior, and optimal use of the report. It does not predict human performance.

\textbf{Stage B: communication comparison.} Run interactive tasks with matched final rights and evidence access. Human mediation can change delay and adaptation. Measure these effects rather than assuming the replay distribution remains valid.

\textbf{Stage C: rights comparison.} Vary whether corrective choices can actually change or stop execution, while holding reports fixed where feasible. Report achievable outcomes under each rights arrangement alongside information regret. The previous example shows why either measure alone is inadequate.

\topic{Operational evidence and normative limits}{Correction success is a testable outcome; legitimate control also requires a stated account of authority.}
\citet[PDF pp. 7--9]{santoni2018control} propose two necessary conditions for meaningful human control: tracking relevant human reasons, and tracing actions to humans with appropriate understanding and responsibility. A click or an audit trail is not sufficient for either. These are philosophical design conditions, not an empirical conclusion about our workflow.

\begin{tabularx}{\linewidth}{L{0.30\linewidth}Y}\toprule
Measure in a proposed pilot & Question it addresses \\\midrule
Correction before execution & Did relevant evidence reach a person in time, and did their response change the action? \\
Value/regret and error costs & Was the decision better under the stated objective? Report stakeholder outcomes separately where objectives conflict. \\
False interventions and delay & What useful coordination was disrupted or slowed? \\
Workload and evidence access & Could people examine material that would change their judgment? \\
Distribution of influence & Whose objections affected outcomes, and who could change the objective? \\\bottomrule
\end{tabularx}

\textbf{Hypothesis.} Reports that preserve relevant corrections, together with usable rights, improve correction under changed conditions at an acceptable attention cost. \textbf{Alternative.} Direct logged collaboration with targeted escalation performs equally well or better. An advantage due solely to a new veto would support the veto, not universal mediation.

\textbf{Failure conditions.} Reviewers rubber-stamp, important facts disappear in handoffs, corrections arrive too late, or nominal permissions cannot affect execution. Separate channels may still share errors because models, evidence, or incentives are shared. Human review can also introduce correlated errors.

The next protocol must predefine numerical thresholds, sampling, and analysis from task stakes and pilot variance. The toy payoff gap supplies no empirical sample-size estimate. No participants have been recruited and no organizational comparison has been run. The defensible recommendation now is to test actual correction capacity and keep its causal ingredients separate.

\topic{Reading map and exact versions}{Use the sources for what was actually read.}
\small
\begin{tabularx}{\linewidth}{L{0.34\linewidth}Y}\toprule
Source & Exact coverage used \\\midrule
Lauffer et al. 2023 & Official PMLR PDF pp. 1--7 of 15. Definitions, Theorem 1 proof, algorithm and initial examples; appendix proofs unread. \\
Hammond et al. 2023 & Published PDF pp. 1, 3, 10--12, 15--17 of 53; remaining full text unread. \\
J\o rgensen et al. 2026 & arXiv v1, PDF pp. 1--13 of 31; appendix proofs not checked. \\
Ehrlich et al. 2026 & arXiv v1, PDF pp. 1--10, 24--25, 34--35 of 38; no independent replication. \\
Rosas et al., AI in a vat & arXiv v1, PDF pp. 1--14 of 38 from prior review. \\
Santoni de Sio and van den Hoven 2018 & Published PDF pp. 1, 7--9 of 14; tracking/tracing account. \\
XOR Labs working draft & Supplied HTML dated 26 October 2026, six substantive presentation sections read earlier. No verified public URL or individual byline. Proposal only. \\
Repository calculation & \texttt{verify\_strategic\_reports.py}; all five deterministic report partitions, rational arithmetic. \\\bottomrule
\end{tabularx}
\normalsize

\textbf{Feedback citation correction.} The public title is \emph{Beliefs in feedback: Factored representations of interacting agents in transformers}. The authors are David A. Ehrlich, Manuel Baltieri, Fernando E. Rosas, and Alexander B. Boyd. The 38-page public v1 is distinct from the 40-page anonymous attachment; page references and experimental details have been checked against the public version rather than transferred blindly.

\textbf{Scope of review.} This is a targeted primary-text review. The broader acquisition log preserves unread sources, failed retrievals, and source hashes. Full texts remain in the local excluded cache; the public package hosts our notes and artifacts, not the downloaded literature.
\references
\end{document}
